% Root draft, not input. Independent audit required.
\subsection{A one-row block with fixed external entrances}

\begin{lemma}\label{C:one-row-fixed-entrance-collapse}
Let $S$ be a full strongly connected component of a finite
category $C$, with at least two objects in $S$. Assume that its minimal group is trivial and
it has one minimal row at each object. Fix an object $a$,
a minimal idempotent $e\in S(a,a)$, the unique arrows
$u_i\in S(a,i)e$, and arrows $v_i\in eS(i,a)$.
Put $e_i=u_iv_i$; then $v_iu_i=e$.
Suppose $e_i p=p$ for every arrow $p:d\to i$ with
$d\notin S$. Replacing $S$ by one trivial object gives
a strictly smaller homotopy-equivalent finite category.
In particular the hypothesis holds when every nonempty
external incoming hom-set is a singleton. There is a dual
statement with one minimal column and fixed external exits.
\end{lemma}
\begin{proof}
Write $D=C\setminus S$. Define a category $\bar C$ with
objects $D\sqcup\{s\}$, retain all full $D$ hom-sets,
and put
\[
 \bar C(s,s)=\{1_s\},\qquad
 \bar C(d,s)=eC(d,a),\qquad
 \bar C(s,d)=C(a,d)e.
\]
Composition with arrows of $D$, and composition through
$s$ between two objects of $D$, are the original products
in $C$. They are well typed. A path leaving and then
reentering $S$ would put its intermediate objects in the
same strong component, so no such path occurs. Thus there
is no composite from $s$ through $D$ back to $s$ to define.
The displayed incoming and outgoing sets are fixed by $e$
on their respective sides. This verifies every associativity
equation involving $1_s$; the others are inherited from $C$.

Define $Q:C\to\bar C$ to fix $D$, send $S$ to $s$ and
every arrow of $S$ to $1_s$, and set
\[
 Q(p:d\to i)=v_i p,\qquad
 Q(q:i\to d)=q u_i.
\]
Since $S(a,i)e$ is a singleton, $f u_i=u_j$ for every
$f:i\to j$ in $S$. For an external entrance $p$,
$p=e_i p=u_i v_i p$. Therefore
\[
 v_j f p=v_j f u_i v_i p=v_j u_j v_i p=v_i p,
 \qquad q f u_i=q u_j.
\]
A composite through $S$ is also preserved, because
$(q u_i)(v_i p)=q e_i p=qp$.
These identities, together with associativity for retained
arrows of $D$, check every typed pair and prove that $Q$
is an identity-preserving functor.

There is a section $J:\bar C\to C$ choosing $a$ for $s$,
the literal displayed arrows for the two kinds of cross hom,
and $1_a$ for $1_s$. The definition of composition makes
this a functor and $QJ=\mathrm{id}_{\bar C}$.
The components $u_i:a\to i$ on $S$ and identities on $D$
give a natural transformation $JQ\Rightarrow\mathrm{id}_C$.
The internal equation is $f u_i=u_j$; the entrance equation
is $u_iv_ip=p$; the exit equation is the definition of $Q(q)$.
Hence $BQ$ and $BJ$ are homotopy inverses.
The quotient loses at least one object and never increases
an incident hom-set or the number of internal arrows, so it
strictly decreases the total arrow count. A singleton incoming
hom-set is necessarily fixed by the endomorphism $e_i$.
Taking opposites proves the dual statement.
\end{proof}

% Further root budget work, not yet a numbered assertion:
% For a no-retract proper two-object SCC of size s in a connected k-object
% category, n>=s+2(k-2)+1. If the SCC has both an incoming and outgoing
% component, the contracted support contains a triangle and two incidences
% with the SCC, giving n>=s+2(k-2)+3.
% At k=8,n<=21 this forces s<=8 at an extreme block and s=6 otherwise.
% An extreme G=1 strong pair of size<=8 is deletable: totals(2,2) collapse;
% totals(3,2)/(2,3) exhaust an eight-arrow minimal band plus identities,
% and a one-row (source) or one-column (sink) object is removed by the
% one-sided corner-deletion lemma. In the nonextreme s=6 equality case,
% every external nonempty hom is singleton, so the lemma above applies.
% Combined with D's audited acyclic excess-six theorem and the existing
% large-SCC exclusion, this will exclude every k>=8 candidate through21.
%
% At k=7, the same method appears to close SCC cases as well. An extreme
% block has s<=10. For s9/10, BS is rationally acyclic; the complement of
% at least5 objects cannot afford a component with nonzero H2 if the
% connected excess<=2 H2-vanishing lemma is supplied (D is checking it).
% Nonextreme blocks have s<=8. For band totals(2,2), s6 needs four
% entrance/exit obstructions to prevent deleting both two-element corners,
% but only two extra arrows are available; an internal extra cross arrow
% can cover two obstructions, while an external parallel extra covers one.
% The s7 cases still leave an unobstructed deletion; at s8 all external
% homs are singleton and the lemma above applies. The alternative s8
% band totals(3,2)/(2,3) permit a one-sided corner deletion directly.
